% PNP REVIEW REPORT::REDTEAM::RT-0013::2026-10-11T15:36:00Z
\documentclass[11pt]{article}
\usepackage[margin=0.83in]{geometry}
\usepackage{amsmath,amssymb,amsthm,booktabs,enumitem,xurl,hyperref}
\setlength{\emergencystretch}{3em}
\newtheorem{theorem}{Theorem}
\newtheorem{corollary}{Corollary}
\newcommand{\status}[1]{\textsc{#1}}
\title{RT-0013: Forward tags do not require inverse incidence\\
\large A black-box locality separation and an all-YES hostile control}
\author{UNIVERSAL\_CORE\_REDTEAM}
\date{11 October 2026, 15:36 UTC}

\begin{document}
\maketitle

\section*{Executive verdict}

\begin{center}
\small
\begin{tabular}{p{0.34\linewidth}p{0.18\linewidth}p{0.40\linewidth}}
\toprule
Claim & Status & Verdict \\
\midrule
\texttt{OFFLINE-BILOCAL-INDEX1} & \status{Proved/Exact} & Accepted at its explicit-input comparison-sort scope. \\
\texttt{FORWARD-TAG-INVERSE-SEPARATION1} & \status{Proved/Exact} & In the deterministic permutation-oracle query model, forward tagged literal access costs one query while inverse incidence needs $n-1$ queries in the worst case. \\
\texttt{INVERSE-INCIDENCE-NECESSARY1} & \status{Falsified/Retired} & Full inverse incidence is not necessary for the tagged second leg: the requested literal's tag can be native to its forward address. \\
\texttt{JMV-PORT-LOCALITY1} & \status{Open} & The exact port/occurrence label in the concrete Jahanjou--Miles--Viola reduction has not been shown locally computable or locally impossible. \\
\texttt{LOCAL-COOK-BILOCAL-M33-1} & \status{Open} & No on-demand, all-cost bridge to the Fortnow--Lipton--van Melkebeek--Viglas model is constructed. \\
\bottomrule
\end{tabular}
\end{center}

The new separation is deliberately model-scoped.  It distinguishes forward
generation, inverse incidence, and offline indexing; it is not a lower bound
for arbitrary SAT algorithms.  Indeed the separating formula family is an
all-YES promise family, so a one-bit promise decider may output YES without
materializing any occurrence object.  P versus NP remains \status{Open}.

\section{Exact candidate and quantifiers}

Let $n=9b$.  On each block of nine preimage labels use the three partitions
\[
\begin{split}
P_0&=\{012,345,678\},\\
P_1&=\{036,147,258\},\\
P_2&=\{048,156,237\}.
\end{split}
\]
For every permutation $\pi:[n]\to[n]$, sign $s\in\{+,-\}$, tag
$t\in\{0,1,2\}$ and triple $T\in P_t$, put into $F_\pi$ the monotone
signed clause $s\pi(T)$.  A literal address contains $(s,t,\text{block},
\text{triple},\text{slot})$.  Therefore its variable is returned with one
query to $\pi$ and its tag $t$ is returned with no additional query.

\begin{theorem}[\texttt{FORWARD-TAG-INVERSE-SEPARATION1}]
For every $n=9b$ and permutation $\pi$, $F_\pi$ has $2n$ proper, distinct
signed 3-clauses and each variable occurs exactly once at every tag for each
sign, hence exactly three times positively and three times negatively.
Forward literal and tag access use one $\pi$-query.  On the other hand, for
every deterministic oracle algorithm $B$ that, for every permutation $\pi$
and target $y$, outputs $\pi^{-1}(y)$, there are $\pi,y$ on which $B$ makes
at least $n-1$ oracle queries.
\end{theorem}

\begin{proof}
Each $P_t$ partitions every nine-element block.  Thus every preimage label
appears exactly once at tag $t$, for each sign.  The displayed triples are
proper and distinct; applying a permutation preserves both properties.

For the inverse bound, answer the first $q<n-1$ distinct queries with
distinct values other than a fixed target $y$.  At least two domain
locations remain unqueried.  There are two permutation completions consistent
with the complete transcript, obtained by placing $y$ at either location and
swapping the remaining images.  Their correct inverse answers differ.
Consequently no such transcript determines $\pi^{-1}(y)$.  Querying $n-1$
locations is sufficient because bijectivity determines the last image.
\end{proof}

This is an exact oracle-query statement.  It does not become a RAM/Turing
time lower bound merely by calling $\pi$ a succinct generator.  If the full
permutation table is part of an explicit encoding, its inverse is obtained by
one scan and $\Theta(n)$ records.  If an inverse table is supplied as advice,
the online lower bound disappears.

\section{Why this falsifies the proposed necessity}

The tagged D\"ocker--Darmann second leg audited in
\texttt{MODELBRIDGE-0013} asks, at a requested output literal, which one of
the three sign-specific occurrences it represents.  The construction above
places that tag in the forward address.  It makes inverse incidence expensive
in the oracle model without making the forward tag expensive.  Therefore
\[
\text{``tag a requested forward literal''}
\not\Rightarrow
\text{``recover every parent incidence of a named variable.''}
\]
Any route that charges a full inverse index to the second leg is charging a
strictly stronger interface than the reduction needs.  The smallest repair is
to isolate the still-open first-leg statement:
\begin{quote}
For the concrete Jahanjou--Miles--Viola clause generator, compute the port or
sign-specific occurrence label of the requested output literal in
$\operatorname{polylog} T$ time and $O(\log T)$ workspace, or prove a
model-matched obstruction for that exact function.
\end{quote}
The preceding theorem neither proves nor refutes that concrete statement.

\section{One-bit, semantic and wrapper hostile controls}

The nine-variable base has the satisfying preimage assignment
\[
(0,0,0,1,0,0,1,0,1).
\]
Repeating it independently in each block and transporting it through $\pi$
satisfies $F_\pi$.  Hence every member of this separating family is
satisfiable.  A promise decider can return one bit, YES, without constructing
tags, incidence lists, a compatibility graph, a state exposure, or a
reconstruction object.  This directly defeats every all-algorithm necessity
claim whose quantified input family includes only these instances.

Clause-order shuffles, within-clause reversal and arbitrary variable
renaming preserve the exact occurrence counts.  An explicit compiler can
sort or scan all $6n$ literal records and construct both directions; a
streaming forward accessor need not.  Padding and redundant clauses can make
the external encoding arbitrarily longer without increasing the semantic
decision cost.  Conversely, storing an inverse table raises representation
size while lowering online access cost.  No semantic-state count therefore
implies a runtime lower bound.

The usual tractable controls remain decisive.  2SAT, Horn and bounded-
treewidth instances admit specialized polynomial algorithms; affine/Tseitin
constraints admit Gaussian elimination; independent products can be solved
componentwise.  None must expose this report's inverse object.  A recognizer
can also dispatch these islands before invoking any general solver, and an
equivalent program wrapper may delay, rename or erase internal states while
preserving its one-bit output.

\section{Bounded reproducible audit}

The script
\texttt{experiments/review/forward\_inverse\_locality\_redteam.py} used a
fixed seed and sizes $n\in\{9,18,27,36\}$, with 32 permutations per size.
It checked 128 instances, 5,760 clauses, 17,280 literal records, 40,192
inverse probes, 256 wrapper transformations and 128 DIMACS-like encoding
round trips.  Every instance had three positive and three negative
occurrences per variable, with tag set $\{0,1,2\}$ for each sign.  The
observed worst inverse scan was exactly $n-1$.  It separately exercised
2SAT, Horn, affine rank, bounded-treewidth, independent-product, padding and
redundancy controls.  Runtime was under $0.06$ seconds; no extrapolation from
these finite checks is used in the theorem.

Script SHA-256:
\texttt{1e37b7cbe0884482ad178115d78e424d8aea7988b5faedf3264dd22d171d06e7}.
Raw-result SHA-256:
\texttt{7a7c84dbd93b4cd6f2dada3755bd5e430189dfa4ded107d5dde5b2dd8f6a8274}.

\section{Primary-literature preflight and transfer audit}

The two blocking obstacles were: (i) the JMV theorem gives highly local
\emph{clause-output} access but does not, in the theorem statement, supply
the inverse-incidence or port-label interface now requested; and (ii) an
offline explicit index is not an on-demand local reduction in the FLMV
random-access model.  The primary sources checked were:
\begin{enumerate}[leftmargin=1.5em]
\item Jahanjou, Miles and Viola, \emph{Local reductions},
\url{https://arxiv.org/abs/1311.3171}.  Its advertised circuit outputs the
$i$th clause with highly local output bits.  This is a directly transferable
clause-access theorem, not an inverse-incidence theorem.
\item Fortnow, Lipton, van Melkebeek and Viglas, \emph{Time-space lower
bounds for satisfiability},
\url{https://lance.fortnow.com/papers/files/tst.pdf}.  Its model-matched
transfer charges random access and local reduction work.  This is directly
relevant, but the missing bi-local first leg remains open.
\item Hellman, \emph{A cryptanalytic time-memory trade-off}, IEEE TIT 26(4),
1980.  Precomputation-versus-online inversion is a useful structural analogy.
Cryptographic hardness is not imported here; only the elementary oracle
adversary above is used.
\item Raman, Raman and Rao, \emph{Succinct indexable dictionaries with
applications to encoding $k$-ary trees, prefix sums and multisets}, TALG
3(4), 2007, and P\u{a}tra\c{s}cu,
\emph{Succincter}, FOCS 2008.  Their rank/select indexes support the
\status{Proved/Exact} offline compiler perspective, but do not convert a
succinct clause generator into a free inverse accessor.
\end{enumerate}

\paragraph{Transfer classification.}
The permutation-oracle adversary and explicit scan/index compiler are
\status{Exact Object/Theorem Transfer} at their stated query and explicit-
input models.  Cryptanalytic inversion tradeoffs and data-structure indexes
are \status{Analogy/Candidate} for designing a charged compiler.  Treating
one-way-function hardness, database terminology, or a random oracle as an
unconditional SAT lower bound is a \status{Failed Transfer}.

\section{Barrier and quantifier audit}

\begin{itemize}[leftmargin=1.5em]
\item \textbf{Relativization.}  The new lower bound is explicitly oracle-
relative.  It therefore cannot itself cross the Baker--Gill--Solovay barrier
or separate P from NP.
\item \textbf{Algebrization.}  No algebraic extension is defeated.  The
argument is a transcript ambiguity for permutations.
\item \textbf{Natural Proofs.}  No largeness/usefulness circuit property is
constructed, so no Natural-Proofs escape is claimed.
\item \textbf{Nonuniformity.}  $\Theta(n\log n)$ advice bits encoding
$\pi^{-1}$ eliminate the online query bound.  Uniform preprocessing must be
charged if it is allowed.
\item \textbf{Arbitrary polynomial exponents.}  The separation does not
turn ``polynomial'' into a fixed exponent and does not compare an unknown SAT
machine to the FLMV endpoint.
\item \textbf{Solver-to-protocol transfer.}  A one-bit decider need not
reveal witnesses, ports, ranks, messages, compatible states or proof nodes.
No black-box extraction theorem is supplied.
\item \textbf{SAME\_REALIZATION.}  Offline build time, table space, online
queries and decision work are kept in one realization.  Mixing the offline
table of one implementation with the workspace of the oracle accessor is
forbidden.
\end{itemize}

\section{Post-quantum technique ledger}

The LWE-style proposal to externally extract structure from black-box solver
behavior remains \status{Analogy/Candidate}: this run shows that even an
exact forward local interface need not expose an inverse interface.  Collision
entropy, Schur squares, quadratic relations, Hilbert functions and special
soundness do not match the current port-locality obstacle and are
\status{Failed Transfer} for this run.  Hash/random-oracle arguments remain
engineering analogies only.  As always, hidden algebraic structure may make
an instance easier rather than provide a hardness source.

\section{Corpus reviewed}

The repository tree was surveyed.  The run read root \texttt{main.tex},
\texttt{AGENTS.md}, \texttt{CODEX.md}, the pnp-research skill, every task
specification, current \texttt{main4.tex}/\texttt{main5.tex}, the newest
verified MAIN material, and current coordination/status files.  Materially
relevant reports included \texttt{LB-0011--0013},
\texttt{MODELBRIDGE-0012--0013}, \texttt{TESTS-0001/0007/0008},
\texttt{RT-0011/0012}, and \texttt{SYNTHESIS-0007/0008}.  Historical HEC,
VSC-HEC, G-PPC-HEC, CRS-PPC and CCSPC causal-branch explorations, governance
notes, and cited experiment artifacts were also inspected.  The immediate
predecessor's initially missing publication was diagnosed as a local
write-procedure anchor mismatch; its preserved body was published and read
back before this audit.

\section{Duplication / contradiction audit}

The raw rank query lower bound in \texttt{LB-0013} is not duplicated: that
report studies occurrence rank from a raw clause stream, whereas this report
separates a native forward tag from inverse access under a permutation
oracle.  \texttt{TESTS-0008} already proved explicit offline indexing; the
present report accepts rather than reproves it and rejects only an unwarranted
necessity inference.  No verified report was found proving
\texttt{JMV-PORT-LOCALITY1}.  The all-YES control is consistent with earlier
Horn, affine, matching and bounded-treewidth counterexamples to universal
resource exposure.

\section{What changed relative to root MAIN}

Root MAIN correctly says P versus NP remains open and identifies the missing
arbitrary-algorithm transfer.  This run makes one local refinement: inverse
incidence must not be listed as necessary for the tagged second leg.  The
highest-value target is now the exact forward port-label function of the
concrete first-leg local reduction, with preprocessing, random access and
workspace charged.  Root \texttt{main.tex} was not edited because this review
role owns no LSPV region or executive synthesis.

\section{Adopted / rejected / retired prior components}

\begin{itemize}[leftmargin=1.5em]
\item \textbf{Adopted:} the tagged second-leg correctness theorem, raw-stream
rank lower bound at its model, explicit offline bi-local compiler, FLMV
model-matching requirement, SAME\_REALIZATION, and recognizable-island
wrapper controls.
\item \textbf{Rejected:} any implication from explicit offline indexing to
on-demand JMV port locality, and any claim that a SAT decider must expose the
index.
\item \textbf{Retired:} \texttt{INVERSE-INCIDENCE-NECESSARY1} for the tagged
second leg.  Full inverse incidence remains optional useful infrastructure.
\end{itemize}

\section{Stale or overstated claims found}

The only overstatement is prospective rather than a verified theorem:
calling the TESTS-0008 inverse table ``the missing object'' would silently
strengthen the second-leg interface.  The report itself properly labels the
on-demand transfer open.  Coordination and manuscript summaries should use
``sufficient offline infrastructure,'' not ``necessary resource.''

\section{Precise missing universal implication}

To obtain P-vs-NP relevance one still needs, for one fixed uniform
deterministic polynomial-time SAT decider $A$, a solver-independent theorem
which converts $A$ and the exact concrete reduction into one FLMV-matched
same-realization algorithm with charged clause access, forward port labels,
preprocessing, reconstruction, time and workspace.  Alternatively, assuming
$P=NP$, one must derive such an algorithm and contradict a model-matched
unconditional lower bound.  No premise may already assert the absence of a
polynomial SAT algorithm.  No such implication is proved here.

\section*{Final status}

\texttt{FORWARD-TAG-INVERSE-SEPARATION1} is \status{Proved/Exact} only in
the declared oracle-query model.  The proposed necessity of full inverse
incidence is \status{Falsified/Retired}.  The JMV port-locality target, the
charged bi-local Cook reduction, the arbitrary-algorithm transfer, and the
P-versus-NP problem remain \status{Open}.

% END_OF_VERIFIED_REPORT::REDTEAM::RT-0013
\end{document}
